
\documentclass{article}
\usepackage{colortbl}
\usepackage{makecell}
\usepackage{multirow}
\usepackage{supertabular}

\begin{document}

\newcounter{utterance}

\centering \large Interaction Transcript for game `bbh', experiment `web\_of\_lies', episode 206 with BSU{-}SLIM\_\_Qwen3.5{-}9B{-}Base\_\_playpen{-}prm{-}9b{-}best\_of\_n.
\vspace{24pt}

{ \footnotesize  \setcounter{utterance}{1}
\setlength{\tabcolsep}{0pt}
\begin{supertabular}{c@{$\;$}|p{.15\linewidth}@{}p{.15\linewidth}p{.15\linewidth}p{.15\linewidth}p{.15\linewidth}p{.15\linewidth}}
    \# & \multicolumn{2}{c}{Player} && \multicolumn{2}{c}{Game Master} \\
    \hline

    \theutterance \stepcounter{utterance}  
    & & & \multicolumn{4}{p{0.6\linewidth}}{
        \cellcolor[rgb]{0.9,0.9,0.9}{
            \makecell[{{p{\linewidth}}}]{
                \texttt{\tiny{[P1$\langle$GM]}}
                \texttt{Evaluate a random boolean function expressed as a word problem.} \\
\\ 
\texttt{Q: Question: Fidel tells the truth. Jerry says Fidel tells the truth. Vina says Jerry tells the truth. Millicent says Vina lies. Raymond says Millicent lies. Does Raymond tell the truth?} \\
\texttt{A: Yes} \\
\\ 
\texttt{Q: Question: Kristian lies. Millie says Kristian lies. Maybelle says Millie tells the truth. Fidel says Maybelle lies. Leda says Fidel lies. Does Leda tell the truth?} \\
\texttt{A: Yes} \\
\\ 
\texttt{Q: Question: Kristian tells the truth. Michaela says Kristian lies. Raymond says Michaela tells the truth. Osvaldo says Raymond tells the truth. Jamey says Osvaldo tells the truth. Does Jamey tell the truth?} \\
\texttt{A: No} \\
\\ 
\texttt{Q: Question: Vernell lies. Conception says Vernell lies. Kandi says Conception lies. Michael says Kandi tells the truth. Osvaldo says Michael tells the truth. Does Osvaldo tell the truth?} \\
\texttt{A:} \\
            }
        }
    }
     \\ \\

    \theutterance \stepcounter{utterance}  
    & \multicolumn{4}{p{0.6\linewidth}}{
        \cellcolor[rgb]{0.561,0.737,0.561}{
            \makecell[{{p{\linewidth}}}]{
                \texttt{\tiny{[P1$\rangle$GM]}}
                \texttt{Let's break down the chain of logic step{-}by{-}step to determine if Osvaldo tells the truth.} \\
\\ 
\texttt{1.  **Vernell**: The premise states that **Vernell lies**.} \\
\texttt{    *   Truth value: False.} \\
\\ 
\texttt{2.  **Conception**: Conception says "Vernell lies."} \\
\texttt{    *   Since Vernell actually lies (False), the statement "Vernell lies" is **True**.} \\
\texttt{    *   Since Conception made a true statement, **Conception tells the truth**.} \\
\texttt{    *   Truth value: True.} \\
\\ 
\texttt{3.  **Kandi**: Kandi says "Conception lies."} \\
\texttt{    *   Since Conception actually tells the truth (True), the statement "Conception lies" is **False**.} \\
\texttt{    *   Since Kandi made a false statement, **Kandi lies**.} \\
\texttt{    *   Truth value: False.} \\
\\ 
\texttt{4.  **Michael**: Michael says "Kandi tells the truth."} \\
\texttt{    *   Since Kandi actually lies (False), the statement "Kandi tells the truth" is **False**.} \\
\texttt{    *   Since Michael made a false statement, **Michael lies**.} \\
\texttt{    *   Truth value: False.} \\
\\ 
\texttt{5.  **Osvaldo**: Osvaldo says "Michael tells the truth."} \\
\texttt{    *   Since Michael actually lies (False), the statement "Michael tells the truth" is **False**.} \\
\texttt{    *   Since Osvaldo made a false statement, **Osvaldo lies**.} \\
\\ 
\texttt{**Conclusion:**} \\
\texttt{Osvaldo does **not** tell the truth.} \\
\\ 
\texttt{A: **No**} \\
            }
        }
    }
    & & \\ \\

    \theutterance \stepcounter{utterance}  
    & & & \multicolumn{2}{p{0.3\linewidth}}{
        \cellcolor[rgb]{0.9,0.9,0.9}{
            \makecell[{{p{\linewidth}}}]{
                \texttt{\tiny{[GM$|$GM]}}
                \texttt{Target: No} \\
            }
        }
    }
    & & \\ \\

    \theutterance \stepcounter{utterance}  
    & & & \multicolumn{2}{p{0.3\linewidth}}{
        \cellcolor[rgb]{0.9,0.9,0.9}{
            \makecell[{{p{\linewidth}}}]{
                \texttt{\tiny{[GM$|$GM]}}
                \texttt{game\_result = LOSE} \\
            }
        }
    }
    & & \\ \\

\end{supertabular}
}

\end{document}
